\section{Standing setup}\label{a:sec:setup}

Let \(Z_0=1\), and let each individual independently leave two offspring with
probability \(p\) and none with probability \(1-p\).  With
\(S_n=\mathbb P(Z_n>0)\), the offspring PGF
\(\phi(z)=pz^2+(1-p)\) gives
\begin{equation}
 S_{n+1}=2pS_n-pS_n^2,
 \qquad S_0=1.
 \label{a:eq:SnIter}
\end{equation}
Throughout the chapter
\begin{equation}
 \varepsilon=1-2p,
 \qquad r=2p,
 \qquad 0\le p<\tfrac12.
 \label{a:eq:epsdef}
\end{equation}
The endpoint value \(A(0)=\tfrac12\) is understood by continuity from
\(p\downarrow0\).

The Kolmogorov constant is
\begin{equation}
 A(p)=\lim_{n\to\infty}\frac{S_n}{(2p)^n}.
 \label{a:eq:setup-A}
\end{equation}
Since \(\mathbb E Z_n=(2p)^n\), conditioning on the non-zero state gives
\begin{equation}
 \mathbb E[Z_n\mid Z_n>0]=\frac{(2p)^n}{S_n}
 \longrightarrow\frac1{A(p)}.
 \label{a:eq:qsmean}
\end{equation}
The full derivation of the conditional law and its variance belongs to \ChM; the
recap here fixes only the interface needed below.

For later comparison, the continuous-time linear birth--death process with birth
and death rates \(\lambda\) and \(\mu\), matched through the competing-clock
probability \(p=\lambda/(\lambda+\mu)\), has
\begin{equation}
 A_{\mathrm c}(p)=\frac{1-2p}{1-p}.
 \label{a:eq:Acts}
\end{equation}
Its derivation is also given in \ChM, and the resulting closed form provides the
explicit benchmark for the discrete constant used throughout this chapter.
